\documentclass[11pt,a4paper]{article}

% ------------------------------------------------------------------
% Packages
% ------------------------------------------------------------------
\usepackage[utf8]{inputenc}
\usepackage[T1]{fontenc}
\usepackage{lmodern}
\usepackage[margin=2.5cm]{geometry}
\usepackage{amsmath,amssymb,mathtools}
\usepackage{booktabs, multirow}
\usepackage{graphicx}
\usepackage{xcolor}
\usepackage{listings}
\usepackage{caption}
\usepackage[section]{placeins}
\usepackage[colorlinks=true,linkcolor=blue!60!black,citecolor=blue!60!black,urlcolor=blue!60!black]{hyperref}

% ------------------------------------------------------------------
% Macros
% ------------------------------------------------------------------
\newcommand{\E}{\mathbb{E}}
\newcommand{\ind}{\mathbf{1}}
\newcommand{\code}[1]{\texttt{#1}}

% Images extracted from the notebook by make_figures.py ; a framed placeholder is printed while a file is missing.
\newcommand{\figfile}[2]{%
  \IfFileExists{figures/#1}{\includegraphics[width=#2]{figures/#1}}%
  {\fbox{\parbox[c][4cm][c]{#2}{\centering\small missing figure:\\ \texttt{figures/#1}}}}}

\lstset{basicstyle=\ttfamily\footnotesize, keywordstyle=\color{blue!70!black}, commentstyle=\color{gray},
        stringstyle=\color{teal}, frame=single, breaklines=true, language=Python, showstringspaces=false}

% ------------------------------------------------------------------
\title{Cash and Carry on Crypto Perpetual Futures:\\
Funding, Timing and Allocation between BTC and ETH}
\author{Elias Benguedda}
\date{November 2026}

\begin{document}
\emergencystretch=3em
\maketitle

\begin{abstract}
We study the cash and carry on crypto perpetual futures: long the spot, short the same quantity of the perpetual, so
as to earn the funding rate paid every 8 hours by the longs to the shorts. The data are the public archives of
Binance (no account needed) for BTC and ETH, from January 2020 to October 2026, on the 8-hour funding grid. The
return of the position decomposes exactly into a funding part and a basis part, the latter telescoping over a
holding period. Always invested, the carry earns 12.3\% a year on BTC and 14.6\% on ETH with a volatility below 2\%
and a maximum drawdown of 2\%, but its yield has fallen from more than 30\% in 2021 to a few percent since 2025. The
funding is persistent (autocorrelation 0.8 at 8 hours), yet a threshold on its expected value only adds 0.4 point a
year, and a dynamic allocation between BTC and ETH, which adds 1.7 point to the 50/50 benchmark in sample, does not
beat it out of sample (7.07\% against 7.48\% on 2023--2026). The fast signals earn a large basis gain without fees,
which disappears with 2 basis points of fees per leg. With the spot as collateral of the short perpetual, the position
is only liquidated if the index is multiplied by 18.
\end{abstract}

\clearpage
\tableofcontents
\newpage

% ==================================================================
\section{Introduction}
% ==================================================================
A perpetual future is a future without maturity. To keep its price close to the spot, the exchange makes the two
sides pay each other a \emph{funding rate} at regular intervals: when the perpetual trades above the spot, the longs
pay the shorts, and conversely. Since the demand for leveraged long positions usually exceeds the demand for short
positions, the funding has been positive most of the time, at a level well above the money market rate.

The \emph{cash and carry} captures it: buy the spot, sell the same quantity of perpetual. The position has no
exposure to the price of the asset, and receives the funding as long as it is positive. This report measures what
this simple strategy has earned on Binance since 2020, and whether it can be improved by leaving the market when the
funding is low (a threshold strategy) or by allocating the capital between BTC and ETH according to their funding (a
dynamic allocation). The scripts and the notebook are open source (Section~\ref{sec:implementation}).

% ==================================================================
\section{Perpetual futures and the funding rate}\label{sec:funding}
% ==================================================================

\subsection{Funding}
On Binance USD-M perpetuals, the funding is exchanged every 8 hours, at 00:00, 08:00 and 16:00 UTC. At each funding
time $t_k$, a position of $q$ units of perpetual pays
\begin{equation}
q\, F_{t_k}\, f_k
\end{equation}
to the other side if it is long, and receives it if it is short, $F$ being the mark price of the perpetual and $f_k$
the funding rate. The rate is set from the premium of the perpetual over the index (a basket of spot prices), averaged
over the funding interval:
\begin{equation}
f_k = P_k + \mathrm{clamp}\big(I - P_k,\ -0.05\%,\ 0.05\%\big), \qquad I = 0.01\%,
\end{equation}
where $P_k$ is the average premium index and $I$ an interest rate component of 0.01\% per 8 hours (10.95\% a year).
When the premium is small, $f_k = I$: this is why a third of the payments are exactly 0.01\%. The rate is capped at
$\pm 0.30\%$ per 8 hours for BTC.

\subsection{Cash and carry}
At $t_k$, the strategy invests one unit of capital: it buys $q = 1/S_k$ unit of spot and sells $q$ units of
perpetual. Over $(t_k, t_{k+1}]$, its return is
\begin{equation}
r_{k+1} = \underbrace{\frac{f_{k+1}\, F_{k+1}}{S_k}}_{\text{funding}}
+ \underbrace{\frac{(S_{k+1} - S_k) - (F_{k+1} - F_k)}{S_k}}_{\text{basis}} .
\label{eq:return}
\end{equation}
The price of the asset cancels out: only the funding received and the change of the basis $b = F - S$ remain. Over a
holding period from $t_a$ to $t_b$, without rebalancing, the basis part telescopes to
\begin{equation}
q \sum_{k=a}^{b-1} \big[(S_{k+1} - S_k) - (F_{k+1} - F_k)\big] = q\,(b_a - b_b):
\end{equation}
the strategy gains the basis at the entry and gives back the basis at the exit. A strategy that enters when the
perpetual is expensive and exits when it is cheap earns this difference on top of the funding.

\subsection{Margin}
Both legs are held on the same exchange in cross (portfolio) margin: the spot is the collateral of the short
perpetual, valued with a haircut $h$, and the perpetual requires a maintenance margin $m$ of its notional. With a
quantity $q$ fixed at the entry, the equity and the maintenance margin are
\begin{equation}
E_t = h\, q\, S_t + q\,(F_{\mathrm{entry}} - F_t) + \text{funding received}, \qquad M_t = m\, q\, F_t,
\end{equation}
and the position is liquidated when $E_t < M_t$. A rise of the index creates a loss on the perpetual which is covered
by the spot. Without funding and with $F \approx S = x\, S_{\mathrm{entry}}$, the condition $E_t \ge M_t$ reads
\begin{equation}
h x + 1 - x \ge m x \iff x \le \frac{1}{1 - h + m},
\end{equation}
that is a multiple of 18.5 for $h = 95\%$ and $m = 0.4\%$. The funding received since the entry pushes the limit
further.

% ==================================================================
\section{Data}\label{sec:data}
% ==================================================================
The data come from the public archives of Binance, \code{data.binance.vision}, which require no account and no API
key: spot klines of BTCUSDT and ETHUSDT, klines and premium index of the USD-M perpetuals, and the history of the
funding rates, from January 2020 to October 2026. The prices are taken at each funding time (open of the 8-hour kline
starting at $t_k$), which gives 7{,}461 dates per symbol; 27 dates are missing from the archives (a few days of 2021)
and are skipped. Table~\ref{tab:funding} and Figure~\ref{fig:funding} describe the funding.

\begin{table}[h]
\centering
\small
\begin{tabular}{@{}lrr@{}}
\toprule
 & \textbf{BTCUSDT} & \textbf{ETHUSDT} \\
\midrule
Mean (\% p.a.) & 11.68 & 13.73 \\
Median (\% p.a.) & 9.80 & 10.77 \\
Standard deviation (\% p.a.) & 22.72 & 29.68 \\
Negative payments (\%) & 13.95 & 13.83 \\
Payments at exactly 0.01\% (\%) & 33.56 & 33.39 \\
Maximum / minimum (\% per 8 hours) & $0.300$ / $-0.300$ & $0.375$ / $-0.356$ \\
\midrule
Autocorrelation at 8 hours / 1 day & 0.80 / 0.70 & 0.79 / 0.66 \\
Autocorrelation at 1 week / 30 days & 0.47 / 0.33 & 0.45 / 0.32 \\
\midrule
Funding received in 2020 / 2021 (\%) & 17.2 / 30.7 & 27.5 / 37.4 \\
2022 / 2023 (\%) & 4.1 / 7.9 & 0.8 / 8.2 \\
2024 / 2025 / 2026 to October (\%) & 12.0 / 5.1 / 2.6 & 13.0 / 4.9 / 1.7 \\
\bottomrule
\end{tabular}
\caption{Funding rates of the Binance perpetuals, January 2020 -- October 2026.}
\label{tab:funding}
\end{table}

\begin{figure}[h]
\centering
\figfile{funding.png}{\textwidth}
\caption{Funding rates, annualised, and their 30-day mean.}
\label{fig:funding}
\end{figure}

The funding is positive on average, around 12\% a year, but very unequal across the years: more than 30\% in 2021,
almost nothing in 2022 and since 2025. It is also persistent, which is the condition for a timing strategy to work.
The basis between the perpetual and the spot at the funding times is small: $-1.6$ bp on average for BTC, with a
standard deviation of 5.8 bp and 99\% of the observations between $-8$ and $+15$ bp (Figure~\ref{fig:basis}).

\begin{figure}[h]
\centering
\figfile{basis.png}{0.8\textwidth}
\caption{Distribution of the basis perpetual / spot at the funding times.}
\label{fig:basis}
\end{figure}

% ==================================================================
\section{Benchmark: always invested}\label{sec:benchmark}
% ==================================================================
The benchmark holds the carry at all times, on BTC, on ETH, or half of the capital on each (Table~\ref{tab:benchmark},
Figure~\ref{fig:benchmark}). The returns are compounded every 8 hours; the volatility and the Sharpe ratio are
annualised with $3 \times 365$ periods a year, the Sharpe ratio being computed without risk-free rate.

\begin{table}[h]
\centering
\small
\begin{tabular}{@{}lrrrrrr@{}}
\toprule
 & \textbf{Return} & \textbf{Volatility} & \textbf{Sharpe} & \textbf{Max drawdown} & \textbf{Funding} & \textbf{Basis} \\
 & (\% p.a.) & (\% p.a.) & & (\%) & (\% p.a.) & (\% p.a.) \\
\midrule
Always BTC & 12.33 & 1.49 & 7.8 & $-2.23$ & 11.67 & $-0.03$ \\
Always ETH & 14.56 & 1.66 & 8.2 & $-2.01$ & 13.78 & $-0.18$ \\
Always 50/50 & 13.44 & 1.39 & 9.1 & $-1.35$ & 12.73 & $-0.11$ \\
\bottomrule
\end{tabular}
\caption{Benchmark, January 2020 -- October 2026, without fees.}
\label{tab:benchmark}
\end{table}

\begin{figure}[h]
\centering
\figfile{benchmark.png}{\textwidth}
\caption{Benchmark: equity (log scale) and drawdown.}
\label{fig:benchmark}
\end{figure}

The return is the funding: the basis contributes almost nothing over the whole period, since it telescopes. The
volatility is very low, but the yield is not stable. On BTC, the carry earned 18.5\% in 2020 and 36.1\% in 2021, then
4.2\% in 2022, 8.0\% in 2023, 12.8\% in 2024, 5.3\% in 2025 and 2.6\% from January to October 2026. The drawdowns come
from the short periods of negative funding and from the moves of the basis in market crashes (March 2020, May 2021):
the worst 8-hour return of the 50/50 benchmark is $-1.1\%$.

% ==================================================================
\section{Funding threshold}\label{sec:threshold}
% ==================================================================

\subsection{Signals}
The threshold strategy holds the carry on a symbol only when its expected funding, known at $t_k$, is above an entry
threshold, and keeps it while the expected funding stays above an exit threshold (hysteresis). Four estimators of the
next funding are compared:
\begin{itemize}
\item \code{last}: the funding just paid at $t_k$;
\item \code{mean}: the mean of the last 21 payments (one week);
\item \code{ewm}: an exponentially weighted mean with a half-life of 21 payments;
\item \code{premium}: the Binance formula applied to the premium index at $t_k$.
\end{itemize}

\begin{table}[h]
\centering
\small
\begin{tabular}{@{}lrrrrrrr@{}}
\toprule
\textbf{Strategy} & \textbf{Return} & \textbf{Vol.} & \textbf{Sharpe} & \textbf{Funding} & \textbf{Basis} & \textbf{Invested} & \textbf{Entries} \\
 & (\% p.a.) & (\% p.a.) & & (\% p.a.) & (\% p.a.) & (\%) & \\
\midrule
Always 50/50 & 13.44 & 1.39 & 9.1 & 12.73 & $-0.11$ & 100.0 & 2 \\
Threshold \code{last} & 14.33 & 1.30 & 10.4 & 12.97 & 0.43 & 88.5 & 679 \\
Threshold \code{mean} & 13.51 & 1.35 & 9.4 & 12.71 & $-0.02$ & 87.3 & 39 \\
Threshold \code{ewm} & 13.44 & 1.36 & 9.3 & 12.68 & $-0.06$ & 91.1 & 15 \\
Threshold \code{premium} & 15.77 & 1.22 & 12.0 & 12.29 & 2.37 & 84.3 & 1{,}967 \\
\bottomrule
\end{tabular}
\caption{Threshold strategies, half of the capital on each symbol, entry at 5\% p.a., exit at 0, without fees.}
\label{tab:threshold}
\end{table}

With smoothed signals (\code{mean}, \code{ewm}), the threshold strategy trades a few dozen times in six years and
earns the same return as the benchmark. With the fast signals (\code{last}, \code{premium}), it trades almost every
day and earns a basis gain: it enters when the perpetual is expensive and exits when it is cheap, which the
telescoping of the basis turns into a profit (2.4\% a year with the \code{premium} signal). This gain is made of a few
basis points per trade (Section~\ref{sec:fees}).

\subsection{Choice of the threshold}
Figure~\ref{fig:threshold} and Table~\ref{tab:sweep} move the entry threshold of the \code{mean} strategy, the exit
being at 0 (or at the entry when it is negative).

\begin{figure}[h]
\centering
\figfile{threshold_sweep.png}{\textwidth}
\caption{Threshold strategy (\code{mean} signal) as a function of the entry threshold.}
\label{fig:threshold}
\end{figure}

\begin{table}[h]
\centering
\small
\begin{tabular}{@{}rrrrrr@{}}
\toprule
\textbf{Entry (\% p.a.)} & \textbf{Return (\% p.a.)} & \textbf{Volatility (\% p.a.)} & \textbf{Sharpe} & \textbf{Invested (\%)} & \textbf{Entries} \\
\midrule
$-5$ & 13.72 & 1.37 & 9.43 & 99.6 & 30 \\
0 & 13.84 & 1.36 & 9.58 & 95.7 & 72 \\
2 & 13.75 & 1.36 & 9.54 & 92.7 & 50 \\
5 & 13.51 & 1.35 & 9.42 & 87.3 & 39 \\
8 & 12.89 & 1.34 & 9.07 & 78.5 & 28 \\
11 & 11.58 & 1.33 & 8.31 & 53.9 & 15 \\
15 & 10.99 & 1.31 & 7.99 & 51.8 & 12 \\
20 & 10.26 & 1.30 & 7.53 & 42.6 & 11 \\
30 & 8.91 & 1.28 & 6.73 & 29.8 & 9 \\
\bottomrule
\end{tabular}
\caption{Threshold strategy (\code{mean} signal, exit at 0) as a function of the entry threshold, without fees.}
\label{tab:sweep}
\end{table}

The best thresholds are close to zero: leaving the market only pays when the expected funding is negative. Higher
thresholds leave the market too often and lose more positive funding than the negative periods they avoid, so both
the return and the Sharpe ratio decrease. At its best, the threshold strategy adds 0.4 point a year to the benchmark:
the negative funding periods are short, and their cost is small compared to the funding of the rest of the time.

% ==================================================================
\section{Dynamic allocation between BTC and ETH}\label{sec:dynamic}
% ==================================================================
The dynamic allocation puts all the capital, at each period, on the symbol with the highest expected funding among
those above the entry threshold; the current symbol is only replaced if the other one is better by more than 2\% a
year, to avoid switching back and forth, and the capital stays in cash when no symbol is eligible
(Table~\ref{tab:dynamic}). The variant \code{equal} splits the capital between the eligible symbols.

\begin{table}[h]
\centering
\small
\begin{tabular}{@{}lrrrrrrr@{}}
\toprule
\textbf{Strategy} & \textbf{Return} & \textbf{Vol.} & \textbf{Sharpe} & \textbf{Max DD} & \textbf{Basis} & \textbf{Invested} & \textbf{Entries} \\
 & (\% p.a.) & (\% p.a.) & & (\%) & (\% p.a.) & (\%) & \\
\midrule
Always 50/50 & 13.44 & 1.39 & 9.1 & $-1.35$ & $-0.11$ & 100.0 & 2 \\
Always ETH & 14.56 & 1.66 & 8.2 & $-2.01$ & $-0.18$ & 100.0 & 1 \\
Dynamic \code{last} & 17.96 & 1.68 & 9.9 & $-1.00$ & 1.72 & 87.6 & 1{,}178 \\
Dynamic \code{mean} & 15.14 & 1.55 & 9.1 & $-1.04$ & $-0.04$ & 85.9 & 69 \\
Dynamic \code{ewm} & 15.02 & 1.57 & 8.9 & $-1.04$ & $-0.20$ & 89.9 & 23 \\
Dynamic \code{premium} & 21.97 & 1.62 & 12.3 & $-0.65$ & 6.00 & 83.1 & 2{,}022 \\
Dynamic \code{mean}, \code{equal} & 13.92 & 1.38 & 9.5 & $-1.17$ & $-0.04$ & 87.3 & 61 \\
\bottomrule
\end{tabular}
\caption{Dynamic allocation between BTC and ETH, entry at 5\% p.a., exit at 0, switching margin 2\% p.a., without
fees.}
\label{tab:dynamic}
\end{table}

\begin{figure}[h]
\centering
\figfile{dynamic_equity.png}{\textwidth}
\caption{Equity of the main strategies (log scale), without fees.}
\label{fig:dynamic}
\end{figure}

\begin{figure}[h]
\centering
\figfile{allocation.png}{\textwidth}
\caption{Dynamic allocation (\code{mean} signal): monthly average weights of BTC and ETH.}
\label{fig:allocation}
\end{figure}

On the one-week mean, the allocation holds ETH 53\% of the time, BTC 33\% and cash 14\%, and follows the funding: 87\%
ETH in 2020 and 2024, mostly BTC in 2022 and 2025 (Figure~\ref{fig:allocation}). It earns 15.1\% a year, 1.7 point
more than the 50/50 benchmark, but only 0.6 point more than the always invested ETH position, which had the highest
funding over the period: most of the gain is the choice of the asset with the higher average funding, not the timing.
The \code{last} and \code{premium} signals earn a large basis gain (1.7\% and 6.0\% a year), the
\code{premium} one reaching 22\% a year with a Sharpe ratio of 12, but with about one trade a day.

% ==================================================================
\section{Risk}\label{sec:risk}
% ==================================================================

\subsection{Worst periods}
The worst 8-hour returns of the dynamic allocation (\code{mean} signal) all occur in market crashes, and come mostly
from the basis, the perpetual falling less than the spot: $-1.00\%$ on 13 March 2020 (funding $-0.33\%$, basis
$-0.67\%$), $-0.62\%$ on 21 December 2020 and on 19 May 2021. These losses are temporary when the position is held,
since the basis converges back, but they are realised by a strategy that exits at that time.

\subsection{Margin}
Table~\ref{tab:margin} gives the lowest ratio of equity to maintenance margin (Section~\ref{sec:funding}) over the
holding periods, with a haircut of 5\% on the spot and a maintenance margin of 0.4\%.

\begin{table}[h]
\centering
\small
\begin{tabular}{@{}lrr@{}}
\toprule
\textbf{Position} & \textbf{Minimum equity / maintenance} & \textbf{Date} \\
\midrule
Always BTC (held from January 2020) & 50.3 & 2021-01-08 \\
Always ETH (held from January 2020) & 29.5 & 2021-05-12 \\
Dynamic \code{mean}, BTC & 168.4 & 2026-09-22 \\
Dynamic \code{mean}, ETH & 93.0 & 2021-05-12 \\
\bottomrule
\end{tabular}
\caption{Lowest margin ratio over the holding periods (liquidation below 1).}
\label{tab:margin}
\end{table}

The ETH position held from January 2020 went through a 33-fold rise of ETH until May 2021, beyond the multiple of
18.5 of Section~\ref{sec:funding}: without funding, the loss of the perpetual would have exceeded the value of the
spot collateral. It still stays 29 times above its maintenance margin, because the funding received over 2020--2021,
paid on a notional that grows with the price, covers the loss. A position entered just before a strong rise, without
this accumulated funding, would be closer to the limit, which remains a rise of the index by a factor 18. With the spot as collateral, the risk of
the strategy is not the liquidation, but the exchange that holds both legs (a default or a freeze of the withdrawals),
and the stablecoin in which the funding is received.

% ==================================================================
\section{Trading fees}\label{sec:fees}
% ==================================================================
The base case has no fees. Table~\ref{tab:fees} adds proportional fees paid on each leg at each change of position;
for comparison, the standard Binance taker fees are 10 bp on the spot and 5 bp on the perpetual, the maker fees being
lower.

\begin{table}[h]
\centering
\small
\begin{tabular}{@{}lrrrrr@{}}
\toprule
\textbf{Fee per leg} & \textbf{0 bp} & \textbf{1 bp} & \textbf{2 bp} & \textbf{5 bp} & \textbf{10 bp} \\
\midrule
Always 50/50 & 13.44 & 13.44 & 13.43 & 13.42 & 13.41 \\
Threshold \code{mean} & 13.51 & 13.39 & 13.26 & 12.89 & 12.26 \\
Threshold \code{ewm} & 13.44 & 13.39 & 13.34 & 13.20 & 12.97 \\
Dynamic \code{mean} & 15.14 & 14.68 & 14.23 & 12.86 & 10.61 \\
Dynamic \code{ewm} & 15.02 & 14.87 & 14.72 & 14.27 & 13.51 \\
Dynamic \code{last} & 17.96 & 10.11 & 2.77 & $-16.43$ & $-40.84$ \\
Dynamic \code{premium} & 21.97 & 8.36 & $-3.73$ & $-32.51$ & $-62.69$ \\
\bottomrule
\end{tabular}
\caption{Annual return (\%) as a function of the fee paid on each leg at each change of position.}
\label{tab:fees}
\end{table}

The fast signals do not survive fees: 2 bp per leg are enough to wipe out the \code{last} signal, and the
\code{premium} signal loses money. The basis gain that they capture is smaller than the cost of crossing the
bid-ask spread twice on two legs. The slow strategies lose 0.2 to 2 points a year at 5 bp: with fees, the dynamic
allocation on the exponential mean (14.27\%) still beats the 50/50 benchmark (13.42\%), but not the always invested ETH
position (14.54\%).

% ==================================================================
\section{Out-of-sample test}\label{sec:oos}
% ==================================================================
The estimator and the entry threshold of the dynamic allocation are chosen on 2020--2022 by the Sharpe ratio, with
fees of 5 bp per leg, then the strategy is run on 2023--2026 (Table~\ref{tab:oos}).

\begin{table}[h]
\centering
\small
\resizebox{\textwidth}{!}{%
\begin{tabular}{@{}llrrrr@{}}
\toprule
\textbf{Signal} & \textbf{Entry} & \textbf{In-sample Sharpe} & \textbf{In-sample return} & \textbf{Out-of-sample Sharpe} & \textbf{Out-of-sample return} \\
 & (\% p.a.) & & (\% p.a.) & & (\% p.a.) \\
\midrule
\code{mean} & 0 & 8.56 & 21.95 & 6.38 & 5.72 \\
\code{mean} & 5 & 8.64 & 21.97 & 7.63 & 6.20 \\
\code{mean} & 8 & 8.79 & 22.04 & 8.62 & 6.26 \\
\code{ewm} & 0 & 9.70 & 24.47 & 9.17 & 6.95 \\
\textbf{\code{ewm}} & \textbf{2} & \textbf{9.70} & \textbf{24.45} & \textbf{9.56} & \textbf{7.07} \\
\code{ewm} & 5 & 9.61 & 24.11 & 10.17 & 7.10 \\
\code{ewm} & 11 & 9.59 & 23.75 & 10.36 & 6.80 \\
\midrule
\multicolumn{4}{@{}l}{Always 50/50, out of sample} & 12.80 & 7.48 \\
\multicolumn{4}{@{}l}{Always ETH, out of sample} & & 7.51 \\
\bottomrule
\end{tabular}}
\caption{Dynamic allocation chosen on 2020--2022 (in bold) and tested on 2023--2026, fees of 5 bp per leg.}
\label{tab:oos}
\end{table}

The parameters chosen in sample earn 7.07\% a year out of sample, against 7.48\% for the 50/50 benchmark and 7.51\% for
the always invested ETH position, with a lower Sharpe ratio (9.6 against 12.8). Over 2023--2026, the funding was lower
and more stable than in 2020--2022, and staying invested was the best choice: the timing does not add value out of
sample.

% ==================================================================
\section{Implementation}\label{sec:implementation}
% ==================================================================
\begin{itemize}
\item \code{strategy/binance\_data.py}: download of the Binance public archives (monthly archives, then daily
archives for the current month), with retries, cache in \code{Data/Crypto} and alignment on the 8-hour funding grid
(the spot timestamps are in microseconds since 2025, the funding times a few milliseconds after the hour);
\item \code{strategy/cash\_and\_carry.py}: returns \eqref{eq:return}, signals, positions (always invested, threshold
with hysteresis, dynamic allocation), backtest with optional fees and cash rate, margin ratio and statistics;
\item \code{strategy/cash\_and\_carry.ipynb}: every number and figure of this report;
\item \code{Reports/CashAndCarry/make\_figures.py}: extracts the figures from the executed notebook.
\end{itemize}

\begin{lstlisting}
from binance_data import load_panels
import cash_and_carry as cc

panels = load_panels(("BTCUSDT", "ETHUSDT"), start="2020-01-01")       # public data, no account
signals = pd.concat({s: cc.funding_signal(panels[s], "ewm", window=21) for s in panels}, axis=1)
weights = cc.dynamic_allocation(signals, entry=0.02, exit=0.0, mode="best", switch=0.02)
result = cc.backtest(panels, weights, fees=(0.0005, 0.0005))
print(cc.performance(result, weights))
\end{lstlisting}

The first download takes about two minutes; afterwards the whole notebook runs in about 15 seconds.

% ==================================================================
\section{Limitations and Next Steps}\label{sec:limitations}
% ==================================================================
\begin{itemize}
\item \emph{Yield of the capital.} The capital is in stablecoins and earns nothing here, nor does the cash of the
timing strategies. With a USD risk-free rate of 4 to 5\% since 2023, the excess return of the carry over the money
market has been small, and the comparison with a cash rate is the next step.
\item \emph{Execution.} Prices are taken at the funding times on 8-hour klines; the fees are proportional, without
bid-ask spread, market impact or slippage between the two legs.
\item \emph{Single exchange.} Both legs are on Binance. The funding differs across exchanges (Bybit, OKX, Hyperliquid),
and an allocation between them is a natural extension, with the cost of moving the collateral and the counterparty
risk of each exchange.
\item \emph{Margin.} The haircut and the maintenance margin are constant; in practice they depend on the size of the
position and on the exchange, and the unrealised loss of the perpetual must be settled in stablecoins.
\item \emph{Universe.} Only BTC and ETH: the funding of the altcoins is higher and more volatile, with liquidity and
delisting risks.
\end{itemize}

% ==================================================================
\section{Conclusion}
% ==================================================================
The cash and carry on Binance perpetuals has earned the funding, 12\% a year on BTC and 14\% on ETH since 2020, with a
volatility below 2\%: a yield well above the money market, but concentrated in 2020--2021 and 2024, and down to a few
percent since 2025. The funding is persistent, yet the timing adds little: a threshold on the expected funding adds
0.4 point a year at best, the dynamic allocation between BTC and ETH mostly captures the higher average funding of
ETH, and the strategy chosen on 2020--2022 underperforms the benchmark on 2023--2026. The fast signals earn a large
basis gain which trading fees wipe out. With the spot as collateral, the liquidation risk is negligible, and the main
risks of the strategy are the exchange, the regimes of low funding and the opportunity cost of the capital.

\end{document}
